% Figura: la matriz de pesos de atención.
% Se usa en sec:marco_atencion, después de la Ecuación de la atención escalada por producto punto.
% Los pesos son ilustrativos: muestran la forma de la operación, no un modelo entrenado.

\begin{figure}[H]
    \centering
    \begin{tikzpicture}[x=1mm, y=1mm]

        % Consultas (filas) y claves (columnas).
        \foreach \i in {1, 2, 3, 4} {
            \node[fignodo, minimum height=6mm, minimum width=9mm, inner sep=1pt,
                  draw=figazul!80!black, fill=figazul!12] (q\i) at (0, 30 - \i*8) {\footnotesize $q_{\i}$};
            \node[fignodo, minimum height=6mm, minimum width=6.5mm, inner sep=1pt,
                  draw=fignaranja!70!black, fill=fignaranja!18] (k\i) at (8 + \i*8, 34) {\footnotesize $k_{\i}$};
        }

        % Pesos softmax(QK^T / sqrt(d_k)); cada fila suma uno.
        \foreach \row/\vals in {1/{60,20,10,10}, 2/{15,55,20,10}, 3/{10,25,50,15}, 4/{5,10,25,60}} {
            \foreach \v [count=\col] in \vals {
                \fill[figceleste!\v!white] (4 + \col*8, 26 - \row*8) rectangle (12 + \col*8, 34 - \row*8);
                \node[font=\scriptsize] at (8 + \col*8, 30 - \row*8) {0.\ifnum\v<10 0\fi\v};
            }
        }
        \draw[figtrazo, shift={(12, -6)}] (0, 0) grid[step=8] (32, 32);

        % Salida.
        \foreach \i in {1, 2, 3, 4} {
            \node[fignodo, minimum height=6mm, minimum width=9mm, inner sep=1pt,
                  draw=figverde!70!black, fill=figverde!10, densely dashed] (o\i) at (66, 30 - \i*8)
                  {\footnotesize $o_{\i}$};
            \draw[figflecha] (45, 30 - \i*8) -- (o\i.west);
        }

        \node[font=\footnotesize, text=figtrazo] at (66, 32) {$o_i = \sum_j a_{ij} v_j$};
        \node[font=\footnotesize, text=figtrazo, align=center] at (28, -12)
            {pesos $a_{ij}$ = $\mathrm{softmax}\!\left(q_i k_j^{\top} / \sqrt{d_k}\right)$};

    \end{tikzpicture}
    \caption{Pesos de atención entre las posiciones de una secuencia de cuatro tokens (valores ilustrativos).}
    \label{fig:marco_atencion}
\end{figure}
